\subsection{多项式代数}

设 B 是一个含单位元的交换结合 A-代数， \((x_{i})_{i\in I}\) 是 B 的一族元素；由诸 \(x_{i}\) ( \(i\in I\) ) 与单位元生成的 B 的子代数记为 \(\mathrm{A}[(x_{i})_{i\in I}]_{\mathrm{B}}\) ，或在不会引起混淆时简记为 \(\mathrm{A}[(x_{i})_{i\in I}]\) 。对每个集合 I，代数 \(\operatorname{Libasc}_{\mathbf{A}}(\mathbf{I})\) 因此等于 \(\mathrm{A}[(\mathbf{X}_{i})_{i\in I}]\) （也记作 \(\mathrm{A}[\mathbf{X}_{i}]_{i\in I}\) )；后一种记号的优点在于指明了用来表示未定元的记号，我们在本专著的其余部分通常采用它。 \(\mathrm{A}[(\mathbf{X}_{i})_{i\in I}]\) 的元素称为未定元 \(X_{i}\) ( \(i\in I\) ) 的、以 A 中元素为系数的多项式；按约定，当说“设 \(\mathrm{A}[(\mathbf{X}_{i})_{i\in I}]\) 为多项式代数”时，总把 \(X_{i}\) 理解为未定元。对 I 的每个子集 J，使用上述记号相当于把 \(\operatorname{Libasc}_{\mathbf{A}}(\mathbf{J})\) 与 \(\operatorname{Libasc}_{\mathbf{A}}(\mathbf{I})\) 的由诸 \(X_{i}\) （ \(i\in J\) ）与单位元生成的代数（参见第 7 号，注记 (3)）等同起来。对于 \(I=\{1,2,\ldots,n\}\) ，我们记 \(\mathrm{A}[\mathbf{X}_{1},\mathbf{X}_{2},\ldots,\mathbf{X}_{n}]\) 以代替 \(\mathrm{A}[(\mathbf{X}_{i})_{i\in I}]\) 。

若 I 与 I' 是两个等势的集合，则代数 \(\operatorname{Libasc}_{\mathbf{A}}(\mathbf{I})\) 与 \(\operatorname{Libasc}_{\mathbf{A}}(\mathbf{I}')\) 同构。 A{[}X{]} 常用来表示对应于一个只含单个元素的未指明指标集的多项式代数，其中 X 表示唯一的未定元；类似地， A{[}X, Y{]} 、 A{[}X, Y, Z{]} 、 \ldots{} 用来表示对应于具有

2, 3, \ldots{} 个元素。注意，根据上述约定，例如 A{[}X{]} 和 A{[}Y{]} 是 A{[}X, Y, Z{]} 的（不同的）子代数，只要 A ≠ \{0\}。

这些元素

\[
\mathrm{X} ^ {\nu} = \prod_ {i \in I} \mathrm{X} _ {i} ^ {\nu (i)},
\]

其中 v 遍历 \(\mathbf{N}^{(1)}\) ，它们构成多项式代数 \(\mathrm{A}[(\mathrm{X}_{i})_{i\in\mathbb{I}}]\) 的一组基。这些元素称为未定元 \(X_{i}\) 中的单项式，而数 \(|v|=\sum_{i\in\mathbb{I}}v(i)\) 称为单项式 \(X^{v}\) 的次数（或总次数）。唯一的 0 次单项式是 \(\mathrm{A}[(\mathrm{X}_{i})_{i\in\mathbb{I}}]\) 的单位元；它常与 A 的单位元 1 等同。 \(\mathrm{A}[(\mathrm{X}_{i})_{i\in\mathbb{I}}]\) 的每个多项式 u 都可以唯一地写成

\[
u = \sum_ {v \in \mathbf {N} ^ {(I)}} \alpha_ {v} X ^ {v}
\]

其中 \(\alpha_{\nu} \in \mathbf{A}\) ；诸元素 \(\alpha_{\nu}\) ——除有限个指标 \(\nu \in \mathbf{N}^{(1)}\) 外均为零——称为 \(u\) 的系数；诸元素 \(\alpha_{\nu} \mathrm{X}^{\nu}\) 称为 \(u\) 的项（元素 \(\alpha_{\nu} \mathrm{X}^{\nu}\) 常称为“ \(\mathrm{X}^{\nu''}\) 中的项”）；特别地，项 \(\alpha_{0} \mathrm{X}^{0}\) （与 \(\alpha_{0} \in \mathbf{A}\) 等同）称为 \(u\) 的常数项。若 \(J\) 是 \(I\) 的子集，则 \(u\) 属于 \(\mathrm{A}[(\mathrm{X}_{i})_{i \in J}]\) 当且仅当 \(\alpha_{\nu} = 0\) 对所有 \(\nu \notin \mathbf{N}^{(J)}\) 成立。由此可得 \(\mathrm{A}[(\mathrm{X}_{i})_{i \in I}]\) 是诸子代数 \(\mathrm{A}[(\mathrm{X}_{i})_{i \in J}]\) 的并，其中 \(J\) 遍历 \(I\) 的有限子集组成的集合。若 \(\alpha_{\nu} = 0\) 对所有 \(|\nu| > n\) 成立，则 \(u\) 称为次数 \(\leqslant n\) 的多项式。当 \(\alpha_{\nu} = 0\) 时，（作为语言上的滥用）称 \(u\) 不含 \(\mathrm{X}^{\nu'}\) 中的项；特别地，当 \(\alpha_{0} = 0\) 时， \(u\) 称为无常数项的多项式。

对每个非零多项式 \(u = \sum_{v} \alpha_v X^v\) ， \(u\) 的次数（或总次数）是整数 \(|v|\) 中最大的一个，其中 \(v\) 是满足 \(\alpha_v \neq 0\) 的多重指标。

设 F 为 A 上的一个含单位元结合代数， \((x_{i})_{i\in I}\) 为 F 的一族两两可交换的元素。F 的子代数 \(F'\) ——由诸 \(x_{i}\) 与单位元生成——是交换的（§1，第 7 号），这使我们可以定义把诸 \(x_{i}\) 代入诸 \(X_{i}\) （在多项式 \(u\in\mathrm{A}[(X_{i})_{i\in I}]\) 中；尽管 F 未必交换）： \(u((x_{i})_{i\in I})\) 是 \(F'\) 的、从而是 F 的元素，而 \(h\colon u\to u((x_{i})_{i\in I})\) 是 \(\mathrm{A}[(X_{i})_{i\in I}]\) 到 F 中的一个同态。同态 h 的核的元素称为族 \((x_{i})\) 在 \(\mathrm{A}[(X_{i})_{i\in I}]\) 中的关系子，也称为诸 \(x_{i}\) 之间的（系数在 A 中的）多项式关系子。同态 h 的像是子代数 \(F'\) ，也记作 \(\mathrm{A}[(x_{i})_{i\in I}]\) （即使 F 不交换时亦然）；若 a 是诸 \(x_{i}\) 之间的多项式关系子的理想，则存在 A-模的正合序列

\[
0 \longrightarrow a \longrightarrow A [ (X _ {i}) _ {i \in I} ] \stackrel {h} {\longrightarrow} A [ (x _ {i}) _ {i \in I} ] \longrightarrow 0.
\]

命题 8. 设 \(\mathbf{A}[(\mathbf{X}_i)_{i\in I}]\) 是一个多项式代数， \(J\) 是 \(I\) 的一个子集， \(K\) 是 \(J\) 在 \(I\) 中的补集。记 \(\mathbf{A}' = \mathbf{A}[(\mathbf{X}_j)_{j\in J}]\) ，并用 \(\mathbf{X}_k'\) ( \(k\in K\) ) 表示多项式代数 \(\operatorname{Libasc}_{\mathbf{A}'}(\mathbf{K}) = \mathbf{A}'[(\mathbf{X}_k')_{k \in \mathbb{K}}]\) 中的未定元，则存在唯一的环同构，它由 \(\mathbf{A}'[(\mathbf{X}_k')_{k \in \mathbb{K}}]\) 映到 \(\mathbf{A}[(\mathbf{X}_i)_{i \in I}]\) 上，在 \(\mathbf{A}'\) 上与恒等映射一致，并把 \(\mathbf{X}_k'\) 映为 \(\mathbf{X}_k\) （对所有 \(k \in \mathbb{K}\) ）。

显然 \(\mathbf{A}[(\mathbf{X}_i)_{i\in I}]\) 是一个 \(\mathbf{A}'\) -代数，由诸 \(\mathbf{X}_k\) （对于 \(k\in \mathbf{K}\) ）生成。另一方面，由于诸 \(\mathbf{X}_k\) ( \(k\in \mathbf{K}\) ) 之间的、系数在 \(\mathbf{A}'\) 中的多项式关系子都可以唯一地写成 \(\sum_{\nu}h_{\nu}((\mathbf{X}_j)_{j\in J})\mathbf{X}^{\nu}\) ，其中 \(\nu\) 遍历 \(\mathbf{N}^{(\mathbb{R})}\) 的一个有限子集，且诸 \(h_\nu\) 是 \(\mathbf{A}[(\mathbf{X}_j)_{j\in J}]\) 的元素，故诸 \(h_\nu\) 必是诸 \(\mathbf{X}_j\) 之间的、系数在 \(\mathbf{A}\) 中的多项式关系子，从而全为零，这就证明了该命题。

命题 8 中所述的同构常被用来把 \(\mathbf{A}[(\mathbf{X}_i)_{i\in I}]\) 的元素与系数在 \(\mathbf{A}' = \mathbf{A}[(\mathbf{X}_j)_{j\in J}]\) 中的多项式等同起来。若 \(u\) 是一个 \(\neq 0\) 的元素（属于 \(\mathbf{A}[(\mathbf{X}_i)_{i\in I}]\) ），则它作为 \(\mathbf{A}'[(\mathbf{X}_k)_{k\in K}]\) 的元素而考虑的总次数也称为它关于诸 \(\mathbf{X}_i\) ——指标为 \(i\in K\) ——的次数。

注记。设 I 与 J 是两个集合， \((\mathrm{P}_{j})_{j\in J}\) 是 \(\mathbf{Z}[(\mathrm{X}_{i})_{i\in I}]\) 的元素的一个族；若 Q 是 \(\mathbf{Z}[(\mathrm{X}_{j})_{j\in J}]\) 的一个使得 \(\mathbf{Q}((\mathrm{P}_{j})_{j\in J})=0\) 的元素，则对环 B 中两两可交换元素的每个族 \((b_{i})_{i\in I}\)

\[
\mathrm{Q} ((\mathrm{P} _ {j} ((b _ {i}) _ {i \in \mathrm{I}})) _ {j \in \mathrm{J}}) = 0.
\]

成立的形如 \(\mathbf{Q}((\mathbf{P}_{j})_{j\in\mathbf{J}})=0\) 的关系有时称为多项式恒等式。例如

\[
\left(\mathrm{X} _ {1} + \mathrm{X} _ {2}\right) ^ {2} - \mathrm{X} _ {1} ^ {2} - \mathrm{X} _ {2} ^ {2} - 2 \mathrm{X} _ {1} \mathrm{X} _ {2} = 0
\]

使得

\[
\begin{array}{r l} \mathrm{Q} & = \mathrm{Y} _ {1} ^ {2} - \mathrm{Y} _ {2}, \qquad \mathrm{P} _ {1} = \mathrm{X} _ {1} + \mathrm{X} _ {2}, \qquad \mathrm{P} _ {2} = \mathrm{X} _ {1} ^ {2} + \mathrm{X} _ {2} ^ {2} + 2 \mathrm{X} _ {1} \mathrm{X} _ {2} \\ & \quad \mathrm{X} _ {1} ^ {n} - \mathrm{X} _ {2} ^ {n} - (\mathrm{X} _ {1} - \mathrm{X} _ {2}) (\mathrm{X} _ {1} ^ {n - 1} + \mathrm{X} _ {1} ^ {n - 2} \mathrm{X} _ {2} + \dots + \mathrm{X} _ {2} ^ {n - 1}) = 0 \end{array}
\]

使得

\[
\begin{array}{c} \mathrm{Q} = \mathrm{Y} _ {1} - \mathrm{Y} _ {2} \mathrm{Y} _ {3}, \qquad \mathrm{P} _ {1} = \mathrm{X} _ {1} ^ {n} - \mathrm{X} _ {2} ^ {n}, \qquad \mathrm{P} _ {2} = \mathrm{X} _ {1} - \mathrm{X} _ {2}, \\ \mathrm{P} _ {3} = \mathrm{X} _ {1} ^ {n - 1} + \mathrm{X} _ {1} ^ {n - 2} \mathrm{X} _ {2} + \dots + \mathrm{X} _ {2} ^ {n - 1} \end{array}
\]

是多项式恒等式。

